\section{Two logarithmic-width binary substitutions}
\label{p2:sec:binary-coding}

The circle words use tags to prevent matches between different children.
We now replace each tagged letter by a binary block. The needed property is
a lower bound for the deficit of every pair of equal-length words, with the
same block-width factor as the elementary upper bound. This makes the loss
in distortion absolute, even as the alphabet grows.

\begin{theorem}[Two binary substitutions]\label{code:two-proofs}
Let $\mathcal A$ be a finite alphabet with $A\geq2$ letters. Each of the two
constructions below gives an integer width $w\geq1$ with $w=O(\log(2A))$, an absolute constant
$a>0$, and a map $c:\mathcal A\to\{0,1\}^w$ such that, simultaneously for
every integer $L\geq0$ and all $x,y\in\mathcal A^L$,
\begin{equation}\label{p2:eq:log-code-deficit}
 a w r_L(x,y)\leq r_{wL}(c(x),c(y))\leq w r_L(x,y).
\end{equation}
Here $c$ extends to words by concatenation. For arbitrary finite words,
\begin{equation}\label{p2:eq:log-code-upper}
 \ED(c(x),c(y))\leq w\ED(x,y).
\end{equation}
The code and its width depend only on $\mathcal A$, not on $L$.
\end{theorem}

The two constructions need not give the same code, width, or constant.
Either substitution may be used in either circle construction.
Their lower bounds imply $(a/2)w\ED(x,y)\leq\ED(c(x),c(y))$ at equal
length, by Lemma~\ref{pre:metric}. The lower assertion is made only at
equal input lengths; the upper assertion holds for arbitrary lengths.
The case $L=0$ is immediate. For the upper deficit bound, concatenate the
code blocks of a common subsequence. For~\eqref{p2:eq:log-code-upper},
simulate each letter insertion, deletion, or substitution by at most $w$
bit edits. It remains to prove the lower deficit bound twice.

\subsection{What an arbitrary bit matching tells us}
Both proofs use the same elementary loss accounting. Fix an increasing
matching between two encoded words of common length $wL$, numbering positions
inside each block from $1$ to $w$, and let $E$ be
its number of unmatched positions on either side. In source block $i$, let
$U_i$ count unmatched positions. If the block has any matches, its target
span is the closed interval from its first partner to its last partner;
let $V_i$ count the globally unmatched target positions in that interval.
For a block without matches put $V_i=0$.

The target spans of different source blocks are disjoint, since their
matches occur in strict order. Also, a matched target position inside one
span has its partner in that source block, by monotonicity. Consequently
\begin{equation}\label{p2:eq:code-loss-sums}
 \sum_iU_i=E,\qquad \sum_iV_i\leq E,
 \qquad |\text{target span of }i|=w-U_i+V_i
\end{equation}
when the span is nonempty. For a threshold $0<\theta<1$, call block $i$
good if $U_i,V_i\leq\theta w$. At most $2E/(\theta w)$ blocks fail this
condition. A good block has a nonempty target span of length at most
$(1+\theta)w<2w$, so its matches meet at most three target blocks.

For the probabilistic choices below, write
$H_2(t)=-t\log_2t-(1-t)\log_2(1-t)$ for binary entropy. We use
\begin{equation}\label{p2:eq:entropy}
 \sum_{j\leq t m}\binom mj\leq2^{H_2(t)m}\qquad(0<t<1/2).
\end{equation}
Indeed, with $z=t/(1-t)<1$,
$(1+z)^m\geq z^{tm}\sum_{j\leq tm}\binom mj$, which rearranges to the
display. The indices in the sum are integers.

\subsection{Different-label intervals and a common offset}
\label{p2:sec:code-offset}
The first code separates sufficiently long intervals of differently
labelled codewords. This kind of interval separation already appears in
the inner code of Schulman and Zuckerman
\cite[Section~II-C, Condition~2 and Lemma~2]{SchulmanZuckerman1999}.
We prove the particular binary property needed here, then use one common
offset in every source block to recover a matching of letters.

Choose an integer $m$ and independent uniform binary words of length $m$, one for each
letter, and put $u=\lfloor m/8\rfloor$. We seek an absolute $\delta>0$
such that every pair of length-$u$ intervals from differently labelled
codewords has longest common subsequence at most $(1-\delta)u$.

Choose an absolute $0<\alpha<1/2$ with
$\gamma=1-\alpha-2H_2(\alpha)>0$. For a fixed pair of intervals with
different labels, the probability of a common subsequence of length
$u-\lfloor\alpha u\rfloor$ is at most
\[
 \binom{u}{\lfloor\alpha u\rfloor}^{\!2}
       2^{-u+\lfloor\alpha u\rfloor}\leq2^{-\gamma u}.
\]
For fixed index lists the involved bits are independent, giving the
probability factor;~\eqref{p2:eq:entropy} bounds their number. There are
at most $A^2m^2$ interval pairs. Taking $m$ to be a sufficiently large
absolute multiple of $\log(2A)$ makes the total failure probability less
than one. For sufficiently large $u$,
$u-\lfloor\alpha u\rfloor\leq(1-\alpha/2)u$; absence of a subsequence of
the former length therefore permits $\delta=\alpha/2$. Fix a code with this property. It makes no assertion
about two intervals of the same codeword.

Choose a positive absolute $\eta<1$ small enough that, for every permitted
width,
\begin{equation}\label{p2:eq:code-offset-parameters}
 2\eta m<\delta\lfloor m/8\rfloor.
\end{equation}
For instance choose $\eta<\delta/100$ first and then increase the minimum
width. Fix $x,y\in\mathcal A^L$, $L\geq1$, and an optimal encoded matching
with deficit $E$. Apply~\eqref{p2:eq:code-loss-sums} with $w=m$ and call a
source block good at threshold $\eta$. At most $2E/(\eta m)$ blocks are
not good.

Let $I=\{\lceil m/3\rceil,\ldots,\lfloor2m/3\rfloor\}$ be the integer
middle-third positions. We claim that a matched position in a good block whose relative index lies in $I$ lands
in a target block with the same letter label. Let $p$ be its relative source
index and $q$ the relative index of its target partner. A length-$u$
interval from $p$ fits in either direction inside the source block. At
least one of those directions fits from $q$ inside its target block. Choose
that direction on both sides. Put $U=U_i$ and $V=V_i$ for this good source
block, and trim the last $V$ source positions from the interval, measured
away from the anchor. For any remaining matched position at source distance
$k\leq u-1-V$ from the anchor, its partner is at target distance at most
$k+V\leq u-1$. The difference of the two distances is the number of
intervening unmatched target positions minus the corresponding source
number, and the former is at most $V$. Monotonicity puts the partner on the
chosen side. At most $U$ retained source positions are unmatched. Hence the
two intervals share a subsequence of length at least
\[
 u-U-V\geq u-2\eta m>(1-\delta)u.
\]
The interval property forces equal labels. This argument also applies when
$q$ is at an endpoint of the target block.

Choose one offset $p$ uniformly from $I$, using this same $p$ in every source
block. The minimum width ensures
$|I|\geq m/4$. Retain each good source block whose position at offset $p$
is matched, and assign it the index of the target block containing that
partner. If $B(p)$ blocks are omitted, then
\begin{equation}\label{p2:eq:code-offset-omissions}
 \E_pB(p)\leq\frac{2E}{\eta m}+\frac E{|I|}
       \leq\left(\frac2\eta+4\right)\frac Em.
\end{equation}
All retained assignments identify equal labels, and their target indices
are nondecreasing. We must count the duplicates before turning them into a
strictly increasing letter matching.

For two consecutive retained blocks that were not originally adjacent,
charge a duplicate target index to an omitted source block strictly between
them. The open source-index gaps of consecutive retained pairs are disjoint,
so these duplicates number at most $B(p)$.

For originally adjacent blocks $i,i+1$, put $d_i=U_i+U_{i+1}$. Their two
positions at a common offset are exactly $m$ apart. If both are matched,
their target partners are at least $m-d_i$ apart. If those partners lie in
one target block, the first is therefore in its first $d_i$ positions.
A good block $i$ meets at most three target blocks, and its distinct source
positions have distinct partners. This duplicate event can thus occur for
at most $3d_i$ offsets; it is impossible when $d_i=0$. Since
$\sum_{i=1}^{L-1}d_i\leq2E$, the number $D_{\rm adj}(p)$ of these adjacent
duplicates satisfies
\begin{equation}\label{p2:eq:code-offset-duplicates}
 \E_pD_{\rm adj}(p)\leq\frac{6E}{|I|}\leq\frac{24E}{m}.
\end{equation}
Only linearity of expectation is used; duplicate events need not be
independent.

In a nondecreasing list, removing one item for each consecutive equality
leaves one item for every distinct index. The retained assignments therefore
give an increasing equal-letter subsequence after at most
$B(p)+D_{\rm adj}(p)$ further deletions. Thus
$r_L(x,y)\leq2B(p)+D_{\rm adj}(p)$. Taking expectations gives
\[
 r_L(x,y)\leq\left(\frac4\eta+32\right)\frac Em.
\]
This proves the first construction in Theorem~\ref{code:two-proofs} with
$w=m$ and $a=(4/\eta+32)^{-1}$. All constants were chosen before $L$,
and the estimate has no additive error depending on $L$.

\subsection{Self-overlaps and an anchor at equal relative indices}
\label{p2:sec:code-anchor}
The second code controls a different local event. It rules out a long,
nearly complete matching between two codeword intervals unless their labels
agree and one matched pair has the same index relative to its two block
starts. That anchor will put a strict majority of one source block into a
single target block, making the recovered letter assignments distinct.

Choose an absolute $\epsilon>0$ with
\begin{equation}\label{p2:eq:code-anchor-parameters}
 2H_2(\epsilon)<\tfrac18,\qquad \epsilon<\tfrac1{10}.
\end{equation}
We seek binary codewords of width $M$ with this precise property: every
increasing matching of at least $M/4$ pairs between two codewords, omitting
at most $\epsilon M$ positions in each of its two spanned intervals, has
equal labels and contains a pair at equal relative indices. The two
codewords may have the same label, so this condition controls shifted
self-overlaps as well.

Self-matching is central to synchronization strings
\cite[Sections~6.1--6.2 of the full version]{HaeuplerShahrasbi2017};
related local conditions occur in the Hamming-to-edit construction of
Bhattacharya et al.\ \cite[Section~2.2]{BhattacharyaEtAl2026}. The property
just stated is the particular one we need, and the following argument
proves it directly.

Choose all code bits independently and uniformly. For a fixed pair of
labels, there are at most
\[
 M^4\left(\sum_{j\leq\epsilon M}\binom Mj\right)^2
 \leq M^4 2^{2H_2(\epsilon)M}
\]
candidate index-list pairs: choose four span endpoints and the omitted
positions; the remaining increasing lists determine the matching. Counting
choices of different list lengths only enlarges the bound.

For different labels, a fixed matching of $q\geq M/4$ pairs has equality
probability $2^{-q}$. For the same label and no equal-index pair, form a
graph on the $M$ independent code bits with one edge for each required
equality. It has no loops. Every index occurs at most once in either list,
so the degree, counting incidences, is at most two. A repeated unordered
edge would require reversed index pairs and violate strict increase; even
if double edges were allowed, a nontrivial connected component with $v$
vertices and $e$ edges has $v\geq2$, $e\leq v$, and $v-1\geq e/2$.
Equality of all bits in that component imposes $v-1$ independent binary
constraints. The full equality system therefore has rank at least $q/2$,
so its probability is at most $2^{-q/2}\leq2^{-M/8}$.

The probability of any violation of the local property is at most
\[
 A^2M^4 2^{(2H_2(\epsilon)-1/8)M}.
\]
By~\eqref{p2:eq:code-anchor-parameters}, this is less than one for $M$ a
sufficiently large absolute multiple of $\log(A+1)$. Fix such codewords.

Now take an optimal matching of $c(x),c(y)$ with deficit $E$, and use
\eqref{p2:eq:code-loss-sums} with $w=M$ and threshold $\epsilon$. At most
$2E/(\epsilon M)$ source blocks are not good. A good block has at least
$(1-\epsilon)M$ matches in a span meeting at most three target blocks.
One target block receives at least $(1-\epsilon)M/3\geq M/4$ of them.
Restrict to matches between this source and target block. A position
between the first and last restricted matches on either side, if globally
matched, is still in the restricted matching: monotonicity places its
partner between the opposite endpoints. Hence the omissions in the two
restricted spans are bounded by $U_i,V_i$. The local property gives equal
labels and an anchor at equal relative indices.

Put $b=\lfloor\epsilon M\rfloor$. Both losses of the good block are at
most $b$. Measure target indices relative to the fixed start of the chosen
target block, even for matches outside it. Relative to the anchor, the
displacement of any match from equal relative indices is the difference
between the intervening unmatched counts. Its absolute value is at most
$2b$. Every matched source index in
$\{2b+1,\ldots,M-2b\}$ therefore still lands in the chosen target block.
At most $b$ of these $M-4b$ positions are unmatched. The chosen block
receives at least
\[
 M-5b\geq(1-5\epsilon)M>M/2
\]
matches. Two good source blocks cannot choose the same target block, since
their matched target positions are disjoint. Their chosen target indices
are increasing by monotonicity. The good blocks thus give an equal-letter
subsequence, proving
\[
 r_L(x,y)\leq\frac{2E}{\epsilon M}.
\]
This is the second construction with $w=M$ and $a=\epsilon/2$, completing
Theorem~\ref{code:two-proofs}.

\subsection{Relation to the earlier alphabet reduction}
The existence order in Theorem~\ref{code:two-proofs} has a close predecessor.
Theorem~1.1 of the full version of Bhattacharya, Dey, Goldenberg, and Kouck\'y
\cite{BhattacharyaDeyGoldenbergKoucky2024} states an alphabet reduction for
normalized insertion--deletion distance, with each distance divided by the
sum of the lengths of its two arguments.
The construction used to prove that theorem is a fixed block substitution
$C:\mathcal A\to Q^k$ in Section~3.2.1 and Lemma~3.3. With approximation
parameter $1/8$, the target alphabet $Q$ has absolute size $q\geq2$, and
$k=O(\log(2A))$. Their stated normalized bounds, applied to that
construction, give
\begin{equation}\label{code:prior-indel}
 \tfrac78 k\Delta(x,y)\leq\Delta(C(x),C(y))\leq k\Delta(x,y).
\end{equation}
Indeed, the total encoded length is $k(|x|+|y|)$, so the normalizing factor
cancels. When both words are empty, the unnormalized inequalities are
immediate.

For clarity, here is the elementary binary step in this comparison.
Give the $q$ target letters distinct $t=\lceil\log_2q\rceil$ bit labels and
encode a letter $a$ by
$j(a)=11\,0\,a_1\,0\cdots a_t\,0$, of fixed width $v=2t+3$.
The string $11$ occurs exactly at a block start. In any increasing bit
matching, let $e$ count unmatched positions in the two words together.
Then at most $e$ source blocks fail to
be fully matched to consecutive target positions: charge a block either
to one of its unmatched source bits or, if all its bits are matched, to
an unmatched target position skipped between consecutive partners. The
charged gaps of different source blocks are disjoint. The same bound holds
on the target side. A fully and consecutively matched block begins at a
genuine target block start, by the unique $11$, so it matches an entire
equal block that is itself good on the target side. Conversely, every good
target block is paired this way. The good blocks therefore give a letter matching with total loss at most
$2e$. Minimizing over bit matchings proves
\[
 \tfrac12\Delta(u,z)\leq\Delta(j(u),j(z))\leq v\Delta(u,z)
\]
for arbitrary finite words; the upper bound simulates letter edits.

The composite $j\circ C$ is a single binary substitution of width $w=vk$.
Combining this inequality with~\eqref{code:prior-indel} gives
$\Delta(jC(x),jC(y))\geq(7/(16v))w\Delta(x,y)$. At equal lengths this is
the deficit lower factor $7/(16v)$, and for ordinary edit distance it gives
the lower factor $7/(32v)$; bitwise simulation gives the upper factor one.
This deduction uses the earlier stated theorem and its block construction.
The two proofs above establish the stated binary bounds directly and exhibit
distinct ways to recover a letter matching.
